\documentclass[11pt]{article}

% Detailed speaker notes for Othmane Zizi (predictive-modelling block, slides 7-11).
% Build:  tectonic docs/presentation/othmane_speaker_notes.tex
\usepackage[margin=0.9in]{geometry}
\usepackage[T1]{fontenc}
\usepackage{lmodern}
\usepackage{xcolor}
\usepackage{enumitem}
\usepackage{titlesec}
\usepackage{fancyhdr}
\usepackage{amssymb}
\usepackage{mdframed}
\usepackage[hidelinks]{hyperref}

\definecolor{bixired}{HTML}{B91C1C}
\definecolor{ink}{HTML}{111827}
\definecolor{slate}{HTML}{374151}
\definecolor{blue}{HTML}{1D4ED8}
\definecolor{blsoft}{HTML}{EFF6FF}
\definecolor{grn}{HTML}{15803D}
\definecolor{paper}{HTML}{F9FAFB}

\newcommand{\arr}{$\rightarrow$}
\newcommand{\x}{$\times$}
\newcommand{\Rtwo}{$R^2$}
\newcommand{\approxx}{$\approx$}

\setlist[itemize]{leftmargin=1.3em, itemsep=2pt, topsep=2pt}
\titleformat{\section}{\large\bfseries\color{blue}}{}{0em}{}
\titlespacing*{\section}{0pt}{16pt}{4pt}

\pagestyle{fancy}
\fancyhf{}
\renewcommand{\headrulewidth}{0.4pt}
\lhead{\small\color{slate}Othmane Zizi --- predictive modelling}
\rhead{\small\color{slate}The Bixi Crew}
\cfoot{\small\color{slate}\thepage}

% "Say this" block
\newmdenv[
  linecolor=blue, linewidth=0.8pt, backgroundcolor=blsoft,
  roundcorner=4pt, skipabove=5pt, skipbelow=6pt,
  innertopmargin=6pt, innerbottommargin=6pt, innerleftmargin=8pt, innerrightmargin=8pt
]{say}

% "Explain simply" block
\newmdenv[
  linecolor=slate, linewidth=0.8pt, backgroundcolor=paper,
  roundcorner=4pt, skipabove=5pt, skipbelow=6pt,
  innertopmargin=6pt, innerbottommargin=6pt, innerleftmargin=8pt, innerrightmargin=8pt
]{simple}

\newcommand{\timing}[1]{\hfill{\normalsize\itshape\color{slate}(#1)}}
\newcommand{\hit}{\smallskip\noindent\textbf{\color{ink}Hit these:}}
\newcommand{\trans}[1]{\smallskip\noindent\textbf{\color{bixired}Transition:} \emph{#1}}

\begin{document}

% ------------------------------------------------------------------ Title
\begin{center}
\vspace*{0.3cm}
{\color{blue}\Huge\bfseries Othmane Zizi}\\[6pt]
{\color{slate}\Large Detailed speaker notes --- predictive modelling}\\[4pt]
{\color{slate}\large Slides 7--11 \;\textbar\; target \approxx 2:00 (stretch 2:30)}\\[10pt]
\rule{0.6\linewidth}{1pt}
\end{center}

\noindent
\begin{minipage}[t]{0.5\linewidth}
\textbf{Your run of slides}
\begin{enumerate}[leftmargin=1.4em, itemsep=1pt]
  \setcounter{enumi}{6}
  \item Modelling approach --- multi-model search
  \item MLflow tracking + registry
  \item Results
  \item Explainability (SHAP + LIME)
  \item Fairness + four-type drift
\end{enumerate}
\end{minipage}\hfill
\begin{minipage}[t]{0.46\linewidth}
\textbf{Numbers cheat-sheet}
\begin{itemize}[leftmargin=1.2em]
  \item Mean target \approxx \textbf{0.33} trips / 15-min slot
  \item Tuned val RMSE \textbf{0.994} dep / \textbf{0.976} arr
  \item Naive val RMSE 1.040 dep / 1.027 arr (\textbf{\approxx 4\% better})
  \item Test \Rtwo{} \textbf{0.334} dep / \textbf{0.339} arr
  \item \Rtwo{} by tier (dep, test): \textbf{0.018 / 0.066 / 0.394}
  \item MLflow runs: \textbf{73} departure / \textbf{57} arrival
\end{itemize}
\end{minipage}

\medskip
\begin{say}
\textbf{Hand-off in (from Louis):} ``Thanks, Louis. So we have clean, leakage-safe feature tables and a
pipeline to train on. I own the modelling --- how we pick a model, track it, and prove it's both
accurate and responsible.''
\end{say}

% ------------------------------------------------------------------ Slide 7
\section{Slide 7 --- Modelling: a multi-model search, auto-selected \timing{\approxx 25--30s}}
\begin{say}
``We deliberately did \emph{not} hand-pick a model. We run a search. We start with two strong
gradient-boosting baselines --- LightGBM and XGBoost. On top of that we add \textbf{FLAML}, an AutoML
layer that searches model families and configurations, and \textbf{Optuna}, which does Bayesian
hyperparameter tuning. \emph{Every} trial is logged to MLflow. The pipeline then selects the best model
purely by \textbf{validation RMSE} and promotes it automatically. We run this once per target ---
departures and arrivals --- and, independently, both land on the same winner: an Optuna-tuned LightGBM,
`lgbm\_optuna'. So the choice is automatic and reproducible, not a guess.''
\end{say}
\hit
\begin{itemize}
  \item Candidates (LightGBM + XGBoost) \textbf{+ FLAML AutoML + Optuna HPO}; every trial \arr{} MLflow.
  \item Selection is \textbf{by validation RMSE}, automatic; both targets independently pick \texttt{lgbm\_optuna}.
\end{itemize}
\trans{``Because every trial is tracked, we can stand behind that choice --- here's the tracking.''}

% ------------------------------------------------------------------ Slide 8
\section{Slide 8 --- MLflow tracking + registry \timing{\approxx 20--25s}}
\begin{say}
``All of that is fully traceable. We run an MLflow tracking server on EC2 with an S3 artifact store. The
departure experiment has \textbf{73 runs}, arrival \textbf{57} --- every candidate, FLAML and Optuna trial,
with its parameters, metrics and the model itself. The single best run per target is registered in the
model registry and promoted to the \textbf{production} alias --- and that production model is exactly the
artifact our apps load and serve. No retraining drift between what we evaluated and what we ship.''
\end{say}
\hit
\begin{itemize}
  \item Tracking server on \textbf{EC2 + S3}; experiments \texttt{bixi-demand-departure} / \texttt{-arrival}.
  \item Best run \arr{} \textbf{registered} \arr{} \textbf{production} alias = the exact artifact served.
\end{itemize}
\trans{``So how good is that production model? Here are the numbers.''}

% ------------------------------------------------------------------ Slide 9
\section{Slide 9 --- Results \timing{\approxx 30--35s --- your anchor slide}}
\begin{say}
``Two things to read here. First, be honest about the metric: 15-minute, per-station demand is genuinely
noisy --- the average slot is about a third of a trip, and most slots are zero --- so absolute \Rtwo{} is
modest \emph{by design}. What matters operationally is the error size, and that's well under \textbf{one
trip per slot}: RMSE around 1.0, typical error about 0.6.

Second, it's a real model, not a baseline. Against a naive historical-average baseline we cut validation
RMSE by about \textbf{4\%} and add roughly six points of \Rtwo. It's \textbf{stable across two unseen
months} --- May and October 2025 --- so no overfitting. And the predictive power is concentrated exactly
where decisions are made: on high-demand stations the \Rtwo{} is \textbf{0.39}, versus near zero on the
quietest ones.''
\end{say}
\hit
\begin{itemize}
  \item Errors \textbf{$<$ 1 trip / slot} (RMSE \approxx 1.0, MAE \approxx 0.6); mean target only \approxx 0.33.
  \item \textbf{Beats naive} (\approxx 4\% RMSE, +6 pts \Rtwo); \textbf{stable} on two unseen months.
  \item \Rtwo{} by tier (dep, test): low 0.018, medium 0.066, \textbf{high 0.394} --- strong where it counts.
\end{itemize}
\trans{``And we can explain \emph{why} it predicts what it does.''}

% ------------------------------------------------------------------ Slide 10
\section{Slide 10 --- Explainability: SHAP \& LIME \timing{\approxx 20--25s}}
\begin{say}
``We explain both models with SHAP and LIME. The SHAP beeswarms tell a clean, sensible story: most of the
signal comes from the \textbf{2024 historical baselines} --- recent-lag and historical-average demand ---
and from the \textbf{cyclical time-of-day} features. Weather is a real but \emph{secondary} driver. It's
the same picture for departures and arrivals, which is reassuring --- the model is leaning on stable
demand structure, not noise. LIME lets us explain any single prediction on demand, which matters for
operator trust.''
\end{say}
\hit
\begin{itemize}
  \item Top drivers: \texttt{baseline\_prev\_15min}, \texttt{hist\_avg\_demand}, cyclical time; weather secondary.
  \item Same story dep \& arr \arr{} sensible behaviour; \textbf{LIME} for per-prediction explanations.
\end{itemize}
\trans{``Explainability leads straight into responsible AI --- fairness and drift.''}

% ------------------------------------------------------------------ Slide 11
\section{Slide 11 --- Fairness \& four-type drift \timing{\approxx 25--30s}}
\begin{say}
``Two responsible-AI checks. \textbf{Fairness:} we measure error parity --- RMSE and MAE --- across demand
tiers and geography. As the tier numbers showed, accuracy concentrates in high-demand stations, and on the
quietest stations \Rtwo{} is near zero. We don't hide that --- we \emph{surface} it so operators don't
over-trust a quiet-station forecast.

\textbf{Drift:} we monitor all four types with Evidently --- \textbf{feature, target, prediction and
concept} drift. On the October-2025 data you can see feature drift is detected on a majority of columns,
which is exactly the kind of seasonal shift the system is meant to catch. That's the hook for scheduled
monitoring and retraining. That's the modelling story --- over to Rui for what we do with it.''
\end{say}
\hit
\begin{itemize}
  \item Fairness = \textbf{error parity} across tiers/geography; weak on low-demand \arr{} flagged, not hidden.
  \item \textbf{All four} Evidently drift types (feature / target / prediction / concept), shown on Oct-2025.
\end{itemize}

\smallskip
\noindent\textbf{\color{blue}Explain it simply --- the ``12-year-old'' version} (use this if a judge asks
``what \emph{is} drift?'' or you want to teach it on the spot):
\begin{simple}
\textbf{Big idea:} ``Drift just asks one question --- \emph{has the world changed since we taught the
model?} We trained it on 2024; now it's October 2025. Is that a different world? We check four things.''
\begin{itemize}[leftmargin=1.4em, itemsep=3pt]
  \item \textbf{1. Feature drift --- did the \emph{inputs} change?} Put all the 2024 temperatures in one
        pile and all the Oct-2025 temperatures in another, and measure how far apart the two piles sit.
        That gap is the \textbf{KS number}: 0 means the piles are identical, 1 means totally different.
        October is a different, colder month, so \texttt{month} (KS 0.75) and \texttt{temperature} (0.30)
        clearly moved --- \textbf{12 of 15} inputs shifted.
  \item \textbf{2. Target drift --- did the \emph{answer} change?} The answer is the real number of trips
        in each 15 minutes. People ride a little differently in autumn, so the thing we're trying to
        predict also shifted a bit (KS 0.11 --- a small move).
  \item \textbf{3. Prediction drift --- did the model's \emph{guesses} change?} If the inputs move, the
        model's outputs naturally move with them (KS 0.36). That's \emph{expected} --- the model is just
        reacting to a new month, it's not a bug.
  \item \textbf{4. Concept drift --- did the \emph{rule} change?} This is the one that really matters. Even
        when all the numbers move, the underlying \emph{rule} --- ``warm weekday evening downtown
        $\to$ lots of trips'' --- might still be true. We test it by asking: \emph{is the model still as
        accurate now as it was before?} Its \Rtwo{} was 0.327 then and 0.334 now --- basically identical
        --- so the rule still holds. \textbf{No concept drift.}
\end{itemize}
\textbf{Punchline:} ``The first three just say \emph{the world looks different} --- keep watching.
Concept drift would say \emph{the model is now actually getting it wrong} --- time to retrain. Here only
the first three fired and the rule held, so the model is \textbf{still trustworthy}; we'd schedule a
retrain as a precaution, not an emergency. (And some of that input drift is expected \emph{by design} ---
we use 2024 as the baseline for every year --- so we review it by hand rather than retraining on a blind
timer.)''
\end{simple}
\trans{``Over to Rui --- turning these forecasts into a rebalancing plan.''}

% ------------------------------------------------------------------ Q&A
\section{Likely questions (be ready)}
\begin{itemize}
  \item \textbf{``Why is \Rtwo{} so low?''} --- It's a 15-minute, per-station target: near-Poisson, mostly
        zeros, huge per-slot variance, so \Rtwo{} understates skill. We predict the conditional mean well
        (MAE \approxx 0.6 trip). It beats the naive baseline, is stable out-of-sample, and reaches
        \Rtwo{} 0.39 on the busy stations where rebalancing decisions actually happen. Aggregate it to
        hourly/station-day and it looks much higher.
  \item \textbf{``Why LightGBM and not the FLAML pick / XGBoost?''} --- We didn't decide; the search did.
        FLAML and Optuna both explored the space; \texttt{lgbm\_optuna} had the lowest validation RMSE for
        \emph{both} targets, so it was auto-selected and promoted.
  \item \textbf{``Isn't target encoding a leakage risk?''} --- Encoders and the 2024 baselines are fit on
        \textbf{train only} (the baselines leave-one-out), then applied forward to May/Oct-2025. Strict
        temporal split, no peeking. (Sarah's slide 3.)
  \item \textbf{``What happens when drift fires?''} --- Evidently flags it; that's the trigger for a
        scheduled re-run of the same pipeline to retrain and re-register a fresh production model --- the
        roadmap item Rui mentions.
  \item \textbf{``Departures and arrivals --- two models or one?''} --- One shared pipeline, run twice
        (once per target). Same features, same selection logic; their difference is the rebalancing signal.
\end{itemize}

\vfill
\noindent{\small\color{slate}\hrulefill\\[4pt]
\textbf{Timing:} 5 slides in \approxx 2 min is brisk --- Slide 9 (Results) is your anchor; if you're
running long, compress Slides 8 and 10 to one sentence each. Keep the naive-baseline and tier-\Rtwo{}
numbers --- they carry the ``it's a real, useful model'' argument.}

\end{document}
